\chapter{Review of Related Works}

This chapter presents a review of existing literature organized thematically to identify the gaps that this study addresses. The themes cover barangay complaint management systems, NLP and machine learning for text classification, Human-in-the-Loop and adaptive learning systems, and workflow and lifecycle management.

\section{Barangay Complaint Management Systems}
Barangay reporting involves community concerns whose administrative handling is distinct from automatic assessment of narrative text. The Local Government Code provides a general local-governance context, while Davao City Reports and the corresponding City Government Citizen's Charter document tracking references, routing and administrative feedback mechanisms \cite{ra7160,dcr}. These sources establish relevant local design precedents; they do not establish that every barangay currently has the same procedures, or that an AI classifier has already improved staff response times. The present study proposes a barangay-specific text triage and auditable case-management design that would require local approval and user testing before operational claims could be made.

\section{NLP and Machine Learning for Text Classification}
The three-class model converts resident-authored narratives into sparse numerical features using TF--IDF, then estimates a researcher-defined text-suitability class using a linear SVM. The main reported study uses word-unigram TF--IDF with the SVM and a separate Naive Bayes baseline; a character-feature ablation is also reported. Word-plus-character fusion and probability calibration are prospective extensions rather than recorded main-model experiments. Scikit-learn documentation describes the chosen vectorizer, estimator and calibration mechanisms \cite{sklearn,sklearntfidf,sklearncalib}.

Philippine complaint pragmatics and code-switching motivate retaining informative words, nonstandard spelling and context rather than equating short or emotional reporting with Spam or incompleteness. De Leon and Parina document variations in Filipino complaint realization, while Bender and Friedman emphasize documenting the language and speaker contexts represented by NLP datasets \cite{deleon,benderfriedman}. These studies do not verify the four assigned language groups in the generated corpus. Synthetic-data literature further warns that generated examples can introduce patterns that inflate performance under the same generator distribution \cite{li2023,shaib}.

\section[Human-in-the-Loop and Adaptive Learning Systems]{Human-in-the-Loop and Adaptive Learning\protect\\ Systems}
Human-in-the-Loop (HITL) systems incorporate human assessment into model-assisted decisions. Prescriptive annotation guidelines and reviewer disagreement are crucial when class definitions require judgments about narrative completeness \cite{rottger2022,fleisig2024}. Cost-sensitive filtering and deferral literature further motivate human verification of potentially consequential Spam predictions \cite{madras2018,franc2023}.

Feedback collection must not be confused with completed active learning. A process that accumulates administrator-verified labels and periodically refits a candidate classifier is more precisely described as \emph{human-feedback-driven batch retraining}. The stronger active-learning term applies when informative samples are deliberately selected for annotation and that sampling procedure is evaluated. Research finds that active-learning gains depend on dataset, model and sampling design; the present study therefore treats the feedback policy as a proposal to validate \cite{ghose2024,beatty2018}.

\section{Workflow and Lifecycle Management}
Case management is an administrative state-transition process, not a fourth or fifth prediction class. Accepted cases can be Unresolved, Ongoing and Resolved; Overdue, Referred and Reopened are logged events rather than text-classifier outputs. Davao City Reports presents a local precedent for history, status tracking and resident follow-up \cite{dcr}. Republic Act No.~6713 and the Anti-Red Tape Authority's explanation of RA 11032 support traceable communications and context-sensitive service processing rules; these should not be misread as universal guarantees that every incident is fully resolved within one fixed interval \cite{ra6713,ra11032}. A proposed workflow therefore requires responsible-person assignment, logged actions, deadlines derived from applicable procedures, evidence review and a reconsideration mechanism.

\section{Theoretical Framework}
The theoretical framework of this study is anchored on six key theories and models: Natural Language Processing (NLP), Text Classification Theory, Term Frequency--Inverse Document Frequency (TF-IDF), Support Vector Machine (SVM), Human-in-the-Loop (HITL) and Active Learning, and State Transition Theory. Each theory is discussed individually below to establish the academic foundation of the proposed system.

\subsection{Natural Language Processing (NLP)}
Natural Language Processing is a field of artificial intelligence that enables computers to transform unstructured human language into structured and machine-readable data. In this study, NLP serves as the foundational layer that processes raw complaint text submitted by barangay residents, allowing it to be analyzed by machine learning models. Studies of Filipino complaint pragmatics and data documentation motivate preserving expression differences and describing speaker/language representation explicitly \cite{deleon,benderfriedman}.

NLP pipelines may include text normalization, tokenization and optional stemming or stopword filtering. In multilingual complaint data, indiscriminate stopword removal or stemming may erase meaningful Cebuano particles and negation. Accordingly, the study retains a conservative text representation and learns TF-IDF vocabulary within the training folds; the figure below describes the actual experimental feature boundary rather than asserting that multilingual stemming was implemented.

The selected proposed preprocessing workflow is specified separately in Section~\ref{sec:proposed-operational-workflows}, without introducing a second overlapping figure.

\subsection{Text Classification Theory}
Text classification theory explains the process of assigning predefined categories to textual data. It operates by learning patterns from labeled training data and applying these learned patterns to unseen inputs. This involves transforming text into numerical representations and using decision-making models to determine the most appropriate category, ensuring consistent and automated evaluation.

In this study, text classification theory underpins the entire complaint categorization process. The system learns from labeled complaint samples to distinguish between valid, suspicious, and spam complaints, enabling consistent and automated classification without relying on manual judgment.

\subsection{Term Frequency--Inverse Document Frequency (TF-IDF)}
TF-IDF is a feature extraction technique that converts textual data into numerical representations based on word importance. It assigns higher weight to words that are significant within a specific document while reducing the influence of common terms that appear across many documents. TF--IDF provides a transparent sparse numerical representation; its value for this corpus is assessed empirically rather than assumed in advance \cite{sklearntfidf}. In this study, TF-IDF transforms multilingual complaint text into structured numerical vectors, enabling the SVM classifier to process the data effectively.

The experimental word-vector baseline and the proposed combined-feature extension are distinguished explicitly in Section~\ref{sec:proposed-operational-workflows}.

\subsection{Support Vector Machine (SVM)}
Support Vector Machine (SVM) is a supervised machine learning algorithm used for classification tasks. It works by identifying an optimal hyperplane that separates different classes in a high-dimensional feature space. SVM is well-suited for text classification because of its effectiveness with high-dimensional and sparse datasets, which are common in NLP applications. For the chosen representation and dataset, the SVM is compared directly with a Multinomial Naive Bayes baseline; the outcome is limited to the specified generated-data evaluation \cite{sklearn}. In this study, SVM operates on TF-IDF-generated features to classify complaints into valid, suspicious, or spam categories.

\subsection{Human-in-the-Loop (HITL) and Active Learning}
The Human-in-the-Loop (HITL) framework integrates human expertise into the machine learning lifecycle to improve model performance and reliability. It can be combined with active learning when a documented query strategy deliberately selects examples for human labeling; merely storing corrections is insufficient to demonstrate improved accuracy.

Batch active-learning research discusses sample selection and stopping conditions, but its gains cannot be assumed for this three-class corpus without a direct test \cite{beatty2018,ghose2024}. In the planned deployment, when barangay administrators correct system-generated classifications, the workflow would record version-specific corrections for quality review and possible future offline retraining. This governed feedback path could allow the system to adapt to local linguistic variations and behavioral patterns over time.



\subsection{State Transition Theory}
State transition theory provides the foundation for workflow management within the system. It models a process as a finite set of states and defines the transitions between them based on specific conditions or actions. In the context of this study, state transition theory governs the complaint lifecycle, defining the movement of a complaint through Unresolved, Ongoing, and Resolved states. Closure, reopening, referral and overdue flags are recorded as administrative events rather than additional prediction classes. This ensures that every complaint is systematically tracked and that barangay officials can monitor progress at each stage, promoting accountability, transparency, and proper resolution.


\section{Barangay Complaint Domains and Legal Jurisdiction}
The Local Government Code (RA 7160) establishes a broad general-welfare and administrative foundation for barangay-level services, dispute-related procedures and community coordination \cite{ra7160}. This foundation must not be misread as authorizing a barangay to settle every report. Solid waste concerns are informed by RA 9003; disaster hazards by RA 10121; adolescent and youth participation by RA 11768; and accounts of financial account scams by RA 12010 \cite{ra9003,ra10121,ra11768,ra12010}. Reports involving violence against women and their children may be meaningful community reports but require safeguards and appropriate authorities; RA 9262 does not make these matters ordinary conciliable barangay disputes \cite{ra9262}. Therefore, this work defines a Valid label to mean a \emph{comprehensible report of a community concern}, irrespective of whether its outcome requires referral. This is a text-classification definition, not a legal finding that a report is factually true.

\section{Complaint Pragmatics: What People Actually Write}
Complaints are communicative acts, not uniform written incident reports. Study of Filipino complaints in English and Tagalog highlights variation in complaint realization and directness \cite{deleon}. Research on anonymous online complaint discourse provides evidence that negative evaluation, criticism, and emotional expression may be part of complaint writing \cite{quimbo}. These studies motivate variation in the intended message: concise description, direct request, rhetorical question, emotional follow-up, secondhand concern and partially ordered narratives. They do not supply representative national percentages for each form. Accordingly, this project \emph{chooses} experimental writing quotas rather than presenting them as natural reporting prevalence. A short or emotional message can be Valid; the intended meaning, not politeness, determines the classification label.

\section{Cebuano, Filipino, English and Code-Switching}
Residents may use Cebuano/Bisaya, Filipino, English and mixed forms. Porlas \cite{porlas} provides a linguistic reference for grammatical patterns in Philippine language code-switching. The generator therefore uses language-associated seed expressions rather than mere random replacement of translated tokens. Some submissions are inherently language-neutral, such as nonsemantic number strings or links. Language-group values in the released dataset refer to generation allocation; they have not been confirmed by independent language annotation. This distinction is important in subgroup reporting.

\section{Controlled Synthetic Data as an Experimental Instrument}
Synthetic examples can deliberately introduce rare or challenging message forms, vary attributes systematically, and preserve provenance without relying on sensitive resident submissions. Long et al. survey synthetic generation, curation and evaluation \cite{long}; Chim et al. investigate style and distributional comparisons for user-generated synthetic data \cite{chim}. Shaib et al. report recurring syntactic templates in generated text, a direct concern for lexical classifiers \cite{shaib}. Counterfactual augmentation can be used to test decision boundaries, but its benefits are not guaranteed. Kaushik et al. describe counterfactual data construction while Huang et al. document a negative result where counterfactually augmented training did not improve generalization \cite{kaushik,huang}. These findings justify \emph{measuring} improvements on a fixed and independent challenge set rather than asserting that synthetic control automatically yields a better deployed model.

\section{Generalization, Annotation Artifacts and Behavioral Testing}
Ribeiro et al. propose CheckList, an NLP testing methodology that probes capabilities using perturbations and expected behavior rather than relying exclusively on average accuracy \cite{ribeiro}. In this project, priority probes include preserving Valid labels under spelling and politeness changes, preserving Spam labels under more natural promotional wording, and changing the class when a previously ambiguous report acquires sufficient essential information. Failure rates on independently written probes help expose missed decision boundaries. Near-perfect synthetic test performance is informative only within the assumptions and lexical distributions of that dataset.

\section{Evaluation Design and Leakage Prevention}
A correct evaluation partitions examples by underlying incident or generation group before fitting text features. A scikit-learn Pipeline places TF-IDF fitting \emph{within} each training fold. Group-aware cross-validation avoids reusing a related scenario in both the corresponding training and validation folds. A separate validation partition supports model selection, and an untouched synthetic test partition reports performance after configuration is fixed. A separately authored challenge set probes distribution shift. Optimizing against the challenge repeatedly would convert it into another validation set \cite{sklearn}. \texttt{LinearSVC} does not publish per-epoch loss; a separate \texttt{SGDClassifier} hinge-loss diagnostic may illustrate convergence but must be clearly attributed to that auxiliary estimator.

\section{Human-in-the-Loop Governance and Operational Readiness}
Amershi et al. describe the importance of human participation in interactive machine learning \cite{amershi}. For barangay complaints, reviewed corrections must be traceable, vetted for mislabeling and sensitive information, and evaluated offline before promotion to production. The approach treats the machine's prediction as provisional. A deployed tool additionally needs authenticated reviewer access, audit retention rules, complaint referrals, error handling, disaster and safeguarding escalation, workflow checks, performance monitoring, and rollback of problematic model versions. These requirements are independent of training dataset size.
